% Part: history
% Chapter: biographies
% Section: alfred-tarski

\documentclass[../../../include/open-logic-section]{subfiles}

\begin{document}

\olfileid{his}{bio}{tar}

\olsection{آلفرد تارسکی}

\olphoto{tarski-alfred}{Alfred Tarski}

آلفرد تارسکی در 14 ژانویهٔ 1901 در ورشوِ لهستان، که آن زمان بخشی از
امپراتوری روسیه بود، به دنیا آمد. تارسکی که او را «ناپلئونی» توصیف
کرده‌اند، پرهیاهو، پرحرف و پرشور بود. انرژی او اغلب در سخنرانی‌هایش
نمایان می‌شد: یک بار هنگام دورانداختن سیگاری در میانهٔ سخنرانی، سطل
زباله‌ای را آتش زد و پس از آن اجازه نیافت در آن ساختمان سخنرانی کند.

تارسکی از جوانی تشنهٔ دانش بود. گرچه بعدها به دانشجویان می‌گفت منطق را
خوانده چون تنها درسی بوده که در آن نمرهٔ~B گرفته است، سوابق دبیرستانش
نشان می‌دهد در همهٔ درس‌ها، حتی منطق، نمرهٔ A گرفته بود. او از 1918 تا
1924 در دانشگاه ورشو تحصیل کرد. تارسکی ابتدا قصد داشت زیست‌شناسی
بخواند، اما به ریاضیات، فلسفه و منطق علاقه‌مند شد، زیرا دانشگاه مرکز
مکتب منطق و فلسفهٔ ورشو بود. او در سال 1924 زیر نظر
Stanis\l{}aw Le\'{s}niewski دکتری گرفت.

تارسکی پیش از مهاجرت به ایالات متحده در سال 1939، هنگامی که در ورشو
معلم دبیرستان بود، بخشی از مهم‌ترین کارهایش را به پایان رساند. کار او
دربارهٔ پیامد منطقی و صدق منطقی در همین دوره نوشته شد. تارسکی در سال
1939 برای دوره‌ای سخنرانی در ایالات متحده بود. در میانهٔ سفر او آلمان
به لهستان حمله کرد و تارسکی به سبب تبار یهودی‌اش نمی‌توانست بازگردد.
همسر و فرزندانش تا پایان جنگ در لهستان ماندند، اما پس از آن توانستند
به ایالات متحده مهاجرت کنند. تارسکی در هاروارد، کالج شهر نیویورک،
مؤسسهٔ مطالعات پیشرفته در پرینستون، و سرانجام دانشگاه کالیفرنیا در
برکلی تدریس کرد. او آنجا برنامهٔ میان‌رشته‌ای منطق و روش‌شناسی علم را
بنیان گذاشت. تارسکی در 26 اکتبر 1983 و در 82سالگی درگذشت.

\begin{reading}
برای آگاهی بیشتر دربارهٔ زندگی تارسکی، زندگی‌نامهٔ
\emph{Alfred Tarski: Life and Logic} را ببینید \citep{Feferman2004}.
آثار بنیادین تارسکی دربارهٔ پیامد منطقی و صدق در ترجمهٔ انگلیسی در
\citep{Tarski1983} در دسترس‌اند. همهٔ آثار اصلی تارسکی در مجموعه‌ای
چهارجلدی گردآوری شده‌اند \citep{Tarski1981}.
\end{reading}

\end{document}
